\documentclass[12pt,a4paper]{article}
\usepackage[utf8]{inputenc}
\usepackage[T1]{fontenc}
\usepackage[brazil]{babel}
\usepackage{geometry}
\usepackage{graphicx}
\usepackage{amsmath}
\usepackage{booktabs}
\usepackage{hyperref}
\usepackage{titleps}
\usepackage{placeins} % Para controlar a posição das figuras
\usepackage{caption} % Melhor controle de legendas
\usepackage[most]{tcolorbox}
\usepackage{mdframed}   % para o box

% Definindo um estilo para o box
\newmdenv[
  linecolor=black,
  linewidth=1pt,
  roundcorner=5pt,
  innertopmargin=10pt,
  innerbottommargin=10pt,
  innerleftmargin=10pt,
  innerrightmargin=10pt,
  frametitlebackgroundcolor=gray!20,
  frametitlefont=\normalfont\bfseries,
  frametitle={Equivalências usando apenas NAND}
]{myformula}
\tcbset{
  myformula/.style={
    colback=blue!5!white,
    colframe=blue!75!black,
    fonttitle=\bfseries,
    title=#1,
    boxrule=0.5mm,
    arc=2mm,
    top=2mm,
    bottom=2mm,
    left=2mm,
    right=2mm
  }
}


\geometry{
    a4paper,
    left=25mm,
    right=25mm,
    top=20mm,
    bottom=20mm
}

% Configuração da página de título
\newcommand{\customtitlepage}{
    \begin{titlepage}
        \vspace*{-0.4cm}
        \begin{figure}
            \centering
            \includegraphics[width=0.5\linewidth]{imagens/logoUFC.png}
        \end{figure}
        \centering
        {\large \textbf{Estatística para engenharia}
        
        
        
        \par}
        \vspace{2cm}
        {\large \textbf{L01 - L03} \par}
        {\large \textbf{Atividades dos TO DO} \par}
        \vspace{5cm}
        
        \vspace{0.5cm}
        \begin{flushleft}
       \large Nome: João Victor de Abreu Silva \large
            
        \end{flushleft}
        \vfill
    \end{titlepage}
}

\begin{document}

\customtitlepage

\section{Slide 1}
\subsection{EXERCISE L1E1}
\begin{table}[h]
    \centering
    \begin{tabular}{|l|l|}
        \hline
        \textbf{Ask} & \textbf{Answer} \\ \hline
        Price of kg of bread in RS & QUALITATIVE \\ \hline
        Family income & QUANTITATIVE \\ \hline
        Salary in RS & QUANTITATIVE \\ \hline
        Income class (I, II, etc) & QUALITATIVE \\ \hline
        Hours of study per day & QUANTITATIVE \\ \hline
        Time for study per day & QUALITATIVE \\ \hline
        Employment level & QUALITATIVE \\ \hline
        Working activity & QUALITATIVE \\ \hline
        Customer satisfaction level & QUALITATIVE \\ \hline
        Eyes colour & QUALITATIVE \\ \hline
        Operating system (of a computer) & QUALITATIVE \\ \hline
        Final grade for the statistic course & QUALITATIVE \\ \hline
        Favoured brand of coffee & QUALITATIVE \\ \hline
        Daily consumption of calories & QUANTITATIVE \\ \hline
        Public debt (of a country) & QUANTITATIVE \\ \hline
        Numb. of components in a family & QUANTITATIVE \\ \hline
        Maximum speed in km/h & QUANTITATIVE \\ \hline
        Shoe size & QUALITATIVE \\ \hline
    \end{tabular}
    \caption{Placeholder caption for ask and answer table for L1E1}
    \label{tab:placeholder_label}
\end{table}

\subsection{EXERCISE L1E2}
Given the following distribution of a generic variable X:\\
\textbf{X -> 2 1 2 1 3 2 0 2 0 1 2 2} \\ \\
QUESTIONS:\\
a)How do you classify the variable?\\
b)Define the number of units and the levels.\\
c)Construct the frequency distribution (absolute, relative and percentage frequency).\\ \\
ANSWERS\\ 
a) QUANTITATIVE DISCRET\\
b) HAVE 12 UNITS AND 4 LEVELS\\
c) the table is on the page below (frequency table)\\ \\ \\ \\ \\
\begin{table}[h]
    \centering
    \begin{tabular}{cccc}
        \toprule
        Value(level) & Absolute frequency ($f_i$) & Relative frequency ($f_i/n$) & Percentual frequency (\%) \\
        \midrule
        0 & 2 & $2/12 = 0.167$ & 16.7\% \\
        1 & 3 & $3/12 = 0.250$ & 25.0\% \\
        2 & 6 & $6/12 = 0.500$ & 50.0\% \\
        3 & 1 & $1/12 = 0.083$ & 8.3\% \\
        \midrule
        Total & 12 & 1.000 & 100\% \\
        \bottomrule
    \end{tabular}
    \caption{Frequency table}
    \label{tab:frequencias}
\end{table}
% Requires: \usepackage{booktabs}


\subsection{EXERCISE L1E3}
The age on a group of 120 men and 80 women was recorded giving the following percentage distribution:\\
\begin{table}[h]
    \centering
    \begin{tabular}{ccc}
        \hline
        \textbf{Age} & \textbf{pMen} & \textbf{pWomen} \\
        \hline
        0 -- 20 & 10 & 20 \\
        20 -- 30 & 10 & 20 \\
        30 -- 50 & 30 & 30 \\
        50 -- 90 & 50 & 30 \\
        \hline
    \end{tabular}
    \caption{Placeholder caption for age distribution by gender}
    \label{tab:age_distribution}
\end{table}

Define:\\
a)The number of statistic units ni younger than 20 years.\\
b)The percentage of statistic units with an age greater or equal to 50 years.\\
c)The number of men that are 30 years old or older.\\


\end{document}
